\BourbakiExerciseGroup

\begin{enumerate}
\def\labelenumi{\arabic{enumi})}
\tightlist
\item
  \begin{enumerate}
  \def\labelenumii{\alph{enumii})}
  \tightlist
  \item
    设 \(\mathbf{K}\) 是特征为 \(p > 0\) 的域。一个有限次的代数扩张 \(\mathbf{F}\) （\(\mathbf{K}\) 上的）被称为例外的，如果它在 \(\mathbf{K}\) 上不是可分的，并且它不包含任何元素 \(x \notin K\) 是 \(p\) -根式的（在 \(K\) 上），换句话说，如果 \(F^p \cap K = K^p\) （V，第151页，习题3）。证明如果 \(F\) 是例外的并且如果 \(E = F_s\) , \(z \in F\) , \(z \notin E\) 且 \(z^p \in E\) ，则对每个 \(t \in K\) 我们有 \(z^p \notin E^p(t)\) （注意，否则我们会有 \(E^p(t) = E^p(t^p)\) 因此 \(t \in F^p\) ).
  \end{enumerate}
\end{enumerate}

\begin{enumerate}
\def\labelenumi{\alph{enumi})}
\setcounter{enumi}{1}
\tightlist
\item
  反之，设E是K的有限次可分代数扩张，并设 \(x \in E\) 使得对于每个 \(t \in K\) 我们有 \(x \notin E^p(t)\) 。设 \(F = E(x^{1/p})\) ，从而 \(E = F_s\) 。证明 F 是 K 的例外扩张。（如果存在 \(y \in F\) 使得 \(y \notin E\) 且 \(y^p = z \in K\) ，证明我们将有 \(z \in E^p(x)\) ；这将蕴含 \(z \in E^p\) ，因此 \(y \in E\) ，这是一个矛盾。）
\end{enumerate}

K 的有限次可分扩张 E 称为强可分的，如果域 \(E^{p}(t)\) （其中 t 遍历 K）的并集不同于 E；因此这等同于说 \(E = F_{s}\) ，其中 F 是 K 的一个不可分例外扩张。K 的强可分扩张不可能是完全域。

\begin{enumerate}
\def\labelenumi{\alph{enumi})}
\setcounter{enumi}{2}
\tightlist
\item
  证明如果 E 是 K 的强可分扩张，则不可能 \([E:E^{p}]=p\) （用反证法论证，表明 \([E:E^{p}]=p\) 将意味着 \(\mathrm{E}^{p}(t)=\mathrm{E}\) 对所有 \(t\in K\) 使得 \(t\notin K^{p}\) ；另一方面注意，如果 \(K\subset E^{p}\) ，则除非 E 是完全的，否则 E 在 K 上不能是强可分的。)
\end{enumerate}

\(d\) ) 证明如果 \([\mathbf{K}:\mathbf{K}^p ] = p\) ，则不存在 \(\mathbf{K}\) 的强可分扩张，因此也不存在 \(\mathbf{K}\) 的例外扩张。

\begin{enumerate}
\def\labelenumi{\arabic{enumi})}
\setcounter{enumi}{1}
\tightlist
\item
  设 \(k\) 是特征为 \(p, a \in k\) 的非完全域，其中 \(a \notin k^p\) ， \(F\) 是有理分式域 \(k(X)\) （ \(k\) 上关于一个未定元的）；令 \(y = X^{p^2} / (X^p + a)\) 与 \(K = k(y)\) .
\end{enumerate}

\begin{enumerate}
\def\labelenumi{\alph{enumi})}
\item
  证明： \([\mathbf{F}:\mathbf{K}] = p^2\) ，从而 \(\mathrm{T}^{p^2} - y\mathrm{T}^p - ay\) 是 \(\mathbf{X}\) 在 \(\mathbf{K}[\mathrm{T}]\) 中的极小多项式；由此推出 \(\mathbf{K}\) 在 \(\mathbf{F}\) 中的相对可分闭包是 \(\mathrm{E} = k(\mathbf{X}^p)\) （参见 V，第 148 页，习题 11）。
\item
  证明 F 是 K 的例外扩张（习题 1）。（用反证法，假设 \(b \in F\) 使得 \(b \notin K\) 且 \(b^{p} \in K\) 。注意 \(\mathbf{K}(b) = k(z)\) ，其中 \(z = r(\mathbf{X}) / s(\mathbf{X})\) ，满足 \(\deg r = p\) , \(\deg s \leqslant p\) 且 r、s 是 \(k[X]\) 中的互素多项式；证明 \(z^{p} \in K\) 并且X在K上的极小多项式为 \(r(T)^{p} - z^{p}s(T)^{p}\) ，相差 K 中的一个因子。利用a)得出结论，我们将有 \(a \in k \cap F^{p} = k^{p}\) ，与假设相反。）
\end{enumerate}

\begin{enumerate}
\def\labelenumi{\arabic{enumi})}
\setcounter{enumi}{2}
\tightlist
\item
  设 K 是特征为 p 的域，p \textgreater{} 0，E是K的有限次扩张， \(E_{r}\) 与 \(E_{s}\) 是包含在 E 中的 K 的最大 p-根式扩张与可分扩张。证明 \(E_{s} \otimes_{K} E_{r}\) 同构于 \(E_{r}\) 在E中的相对可分闭包，并与之等同；若它不等于E，则E是 \(E_{r}\) 的例外扩张，且 \(E_{s} \otimes_{K} E_{r}\) 是 \(E_{r}\) 的强可分扩张。
\end{enumerate}

反之，设E是域L的一个例外扩张，并设K是L的一个子域，使得L在K上是p-根式的，且 \([L:K]<+\infty\) ；则 L 是包含于 E 中的 K 的最大 p-根式扩张。

\begin{enumerate}
\def\labelenumi{\arabic{enumi})}
\setcounter{enumi}{3}
\tightlist
\item
  设 \(\mathbf{K}\) 是一个域， \(\mathbf{E} = \mathbf{K}(x)\) 是 \(\mathbf{K}\) 的一个可分代数扩张，且 \(f\) 是 \(x\) 在 \(\mathbf{K}[\mathbf{X}]\) 中的极小多项式.
\end{enumerate}

\begin{enumerate}
\def\labelenumi{\alph{enumi})}
\item
  设 \(\mathbf{F}\) 是 \(\mathbf{K}\) 的一个扩张；在环 \(\mathbf{F}[\mathbf{X}]\) 中，多项式 \(f\) 分解为乘积 \(f_{1}f_{2}\ldots f_{r}\) ，各因子是不可约、可分且两两互异的多项式。证明代数 \(\mathbf{E} \otimes_{\mathbf{K}} \mathbf{F}\) 同构于域的乘积 \(\mathbf{F}[\mathbf{X}] / (f_j)\) ( \(1 \leqslant j \leqslant r\) ).
\item
  假设 \(\mathbf{F} = \mathbf{K}(y)\) 是 \(\mathbf{K}\) 的可分代数扩张；设 \(g\) 为 \(y\) 在 \(\mathbf{K}[\mathbf{X}]\) 中的极小多项式，并设 \(g = g_1g_2\ldots g_s\) 为 \(g\) 在E{[}X{]}中分解成不可约多项式的乘积。证明 r = s，并且如果 \(m = \deg f\) , \(n = \deg g\) , \(m_j = \deg f_j\) , \(n_j = \deg g_j\) ，则我们在必要时对 \(g_j\) 作置换后，可假设 \(m_j / n_j = m / n\) 对 \(1 \leq j \leq r\) 成立.
\end{enumerate}

\begin{enumerate}
\def\labelenumi{\arabic{enumi})}
\setcounter{enumi}{4}
\tightlist
\item
  设 \(\mathbf{K}\) 是有 \(q\) 个元素的有限域，并令 \(\mathbf{A}\) 为 K-代数 \(\mathbf{K}^n\) 。要使 \(\mathbf{A}\) 由单个元素生成，必要且充分条件是 \(n \leqslant q\) .
\end{enumerate}

